\section{Novikov descent}

We now define the notion of a Novikov descent datum and prove that the category
of Novikov descent data is equivalent to the category of finite projective
modules.

\subsection{Abstract descent data}

\begin{definition}
  \label{def:cosimplicial_ring}
  \lean{Novikov.Descent.Abstract.CosimplicialRing}
  \leanok
  A \emph{cosimplicial ring} (truncated to levels $1$--$3$) consists of
  three commutative rings $R_1, R_2, R_3$ with face maps
  $\pi_1^\ast, \pi_2^\ast \colon R_1 \to R_2$ and
  $\pi_{12}^\ast, \pi_{13}^\ast, \pi_{23}^\ast \colon R_2 \to R_3$,
  defining composites $\rho_i^\ast = \pi_{\cdot}^\ast\pi_{\cdot}^\ast \colon R_1 \to R_3$
  satisfying the six cosimplicial identities.
  For a ring homomorphism $\pi^\ast \colon R \to S$, we write
  $\pi^\ast M = S \otimes_{R, \pi^\ast} M$ for the base change of an $R$-module $M$.
\end{definition}

\begin{definition}
  \label{def:descent_datum}
  \lean{Novikov.Descent.Abstract.DescentDatum}
  \leanok
  \uses{def:cosimplicial_ring}
  A \emph{descent datum} over a cosimplicial ring $C$ is a finite projective
  $R_1$-module $M$ with an $R_2$-linear isomorphism
  $\varphi \colon \pi_1^\ast M \xrightarrow{\sim} \pi_2^\ast M$ satisfying the
  cocycle condition
  $\pi_{23}^\ast\varphi \circ \pi_{12}^\ast\varphi = \pi_{13}^\ast\varphi$.
  Morphisms are $R_1$-linear maps commuting with $\varphi$ after base change;
  they form a category with an internal Hom.
\end{definition}

\begin{definition}
  \label{def:extended_cosimplicial_ring}
  \lean{Novikov.Descent.Abstract.ExtendedCosimplicialRing}
  \leanok
  \uses{def:cosimplicial_ring}
  An \emph{extended cosimplicial ring} adds a ring $R_0$ and
  $\pi_0^\ast \colon R_0 \to R_1$ with
  $\pi_1^\ast\pi_0^\ast = \pi_2^\ast\pi_0^\ast$.
\end{definition}

\begin{definition}
  \label{def:constant_descent_datum}
  \lean{Novikov.Descent.Abstract.constantDescentDatum}
  \leanok
  \uses{def:extended_cosimplicial_ring, def:descent_datum}
  For a finite projective $R_0$-module $M$, the \emph{constant descent datum}
  $\operatorname{Const}(M)$ has underlying $R_1$-module $R_1 \otimes_{R_0} M$,
  with $\varphi$ the composition of associativity isomorphisms
  $\pi_1^\ast(R_1 \otimes_{R_0} M) \cong R_2 \otimes_{R_0} M
   \cong \pi_2^\ast(R_1 \otimes_{R_0} M)$
  via $\pi_1^\ast\pi_0^\ast = \pi_2^\ast\pi_0^\ast$.
\end{definition}

\begin{definition}
  \label{def:constant_descent_functor}
  \lean{Novikov.Descent.Abstract.constantDescentDatumFunctor}
  \leanok
  \uses{def:constant_descent_datum}
  This defines a functor
  $\operatorname{Const} \colon \mathsf{Vect}(R_0) \to \mathsf{Desc}(E)$.
  It preserves internal Hom:
  $\underline{\operatorname{Hom}}(\operatorname{Const}(M), \operatorname{Const}(N))
   \cong \operatorname{Const}(\operatorname{Hom}_{R_0}(M, N))$.
\end{definition}

\begin{theorem}
  \label{thm:constant_descent_fully_faithful}
  \lean{Novikov.Descent.Abstract.constantDescentDatumFunctor_fullyFaithful}
  \leanok
  \uses{def:constant_descent_functor}
  If the natural map $R_0 \to \operatorname{eq}(\pi_1^\ast, \pi_2^\ast)$ is bijective,
  then $\operatorname{Const}$ is fully faithful.
\end{theorem}

\begin{proof}
  \leanok
  \textbf{Faithfulness.}  Bijectivity of $R_0 \to \operatorname{eq}(\pi_1^\ast,\pi_2^\ast)$
  implies $\pi_0^\ast \colon R_0 \to R_1$ is injective, hence
  $p \mapsto 1 \otimes p$ is injective for any finite projective $P$.
  Taking $P = \operatorname{Hom}(M, N)$, two maps $f,g$ equal after
  base change to $R_1$ satisfy $1 \otimes (f-g) = 0$, forcing $f=g$.

  \textbf{Fullness.}  A morphism $F \colon \operatorname{Const}(M) \to
  \operatorname{Const}(N)$ yields via the internal Hom isomorphism an
  element $f \in R_1 \otimes_{R_0} \operatorname{Hom}(M,N)$ fixed by
  $\varphi$.  The equalizer hypothesis plus finite projectivity show
  $f = 1 \otimes f_0$ for some $R_0$-linear map $f_0$, which recovers $F$.
\end{proof}

\subsection{Novikov descent data}

\begin{definition}[Definition~2.9]
  \label{def:novikov_descent_datum}
  \lean{Novikov.Descent.novikovCosimplicialRing}
  \leanok
  \uses{def:cosimplicial_ring, def:novikov_series, def:novikov_substitute}
  Let $A$ be a commutative ring and $\Gamma \subseteq \mathbb{R}$ an additive
  submonoid.  The \emph{Novikov cosimplicial ring} has
  $R_1 = A((t^\Gamma))$,
  $R_2 = A((t^\Gamma, u^\Gamma))$,
  $R_3 = A((t^\Gamma, u^\Gamma, v^\Gamma))$,
  with face maps given by substituting the variables:
  $\pi_1^\ast, \pi_2^\ast \colon R_1 \to R_2$ send $t$ to $t$ or $u$ respectively;
  $\pi_{12}^\ast, \pi_{13}^\ast, \pi_{23}^\ast \colon R_2 \to R_3$ are the three coface
  maps between two and three variables.
  The extended version sets $R_0 = A$ with $\pi_0^\ast \colon A \to A((t^\Gamma))$
  the inclusion as constant series.
  A \emph{Novikov descent datum} is a descent datum over this cosimplicial ring.
\end{definition}

\begin{lemma}
  \label{lem:novikov_equalizer}
  \lean{Novikov.Descent.novikovCosimplicialRing_equalizer,
    Novikov.Descent.novikovExtendedCosimplicialRing_equalizer_bijective}
  \leanok
  \uses{def:novikov_descent_datum}
  The equalizer of
  $\pi_1^\ast, \pi_2^\ast \colon A((t^\Gamma)) \to A((t^\Gamma, u^\Gamma))$
  is exactly the constant series $A$.
  Hence $A \to \operatorname{eq}(\pi_1^\ast, \pi_2^\ast)$ is bijective.
\end{lemma}

\begin{proof}
  \leanok
  If $f(t) \in A((t^\Gamma))$ satisfies $f(t) = f(u)$ in $A((t^\Gamma, u^\Gamma))$,
  evaluate at the exponent $(d, 0)$: the left side gives the
  $d$-coefficient of $f$, while the right side gives the
  $d$-coefficient of the series whose only variable is $u$, hence $0$ for $d \neq 0$.
  Thus all coefficients of $f$ vanish except possibly the constant term, so
  $f \in A$.
\end{proof}

\begin{proposition}[Proposition~2.12]
  \label{prop:vect_to_novikov_descent_fully_faithful}
  \lean{Novikov.Descent.vectToNovikovDescent_fullyFaithful}
  \leanok
  \uses{def:novikov_descent_datum, thm:constant_descent_fully_faithful,
    lem:novikov_equalizer}
  The constant-descent functor
  $\mathsf{Vect}(A) \to \mathsf{NovDD}(A, \Gamma)$ is fully faithful.
\end{proposition}

\begin{proof}
  \leanok
  The equalizer hypothesis of
  Theorem~\ref{thm:constant_descent_fully_faithful}
  is satisfied by Lemma~\ref{lem:novikov_equalizer}.
\end{proof}
